Several directions follow naturally from the present implementation. A
first line of future work concerns completing the symmetry between the
two flight targets. Although the platform already exposes both
Crazyflie-oriented and MAVLink-oriented control paths, a fuller
unification of mission abstractions, telemetry semantics, and safety
behaviors across both environments would strengthen the claim of a truly
portable high-level autonomy layer. This includes deeper integration
with PX4 control modes, more explicit offboard control patterns, richer
use of onboard obstacle maps inside autopilot workflows, and a cleaner
common contract between indoor experimental flights and more operational
autopilot-driven deployments.

On the perception side, future work should extend the current
multi-camera and multi-model logic toward tighter cross-camera
continuity. The present system supports runtime camera switching,
per-camera geometry, and viewpoint-aware projection, but does not yet
provide full track handoff across cameras. A next step would therefore
be to reproject existing track hypotheses across viewpoints and preserve
track identity when a target moves from one sensor to the other. This
would make the dual-camera architecture more than a switching mechanism
and turn it into a more integrated multi-view perception system.

A second major direction concerns model deployment. The current work
argues for compact YOLO-style detectors and open synthetic pre-training
assets, but there remains significant room for improvement in model
adaptation, quantization robustness, and scene-specific specialization.
Future work could evaluate alternative compact detectors, mixed-model
cascades, task-specific switching policies, and better domain adaptation
between synthetic and real aerial imagery. It would also be valuable to
formalize when workloads should remain on the EdgeTPU and when they
should fall back to CPU-side TensorFlow Lite Micro execution, especially
in missions that mix fast detection with lightweight auxiliary models.
In particular, the regulation-aware navigation framing introduced in
this work would benefit from models trained not only for general
detection, but also for explicit recognition of flight-relevant
constraints in inhabited spaces.

The expanded embedded-ML runtime opens a further line of work beyond
classical detector deployment. The current implementation already
exposes clustering, anomaly detection, sequence modeling, reinforcement
learning, and lightweight SLAM primitives, but these remain building
blocks rather than a mature learning framework. Future work could
evaluate how these modules interact in closed-loop autonomy settings,
how they should share memory and telemetry channels with the perception
stack, and which of them are practically robust enough for long-duration
field experiments on an MCU-class platform.

Another important direction is deeper optimization of the memory and
dataflow path. The current implementation already uses task parallelism,
staging buffers, and explicit SDRAM placement, but more aggressive
zero-copy strategies, richer DMA-assisted transfer patterns, or
compiler-assisted layout optimization may further improve throughput and
energy efficiency. The same is true for the host-side interaction with
the EdgeTPU: a richer study of tensor residency, batching patterns, or
asynchronous execution opportunities could reveal additional performance
margins on this class of embedded platform.

The developer experience can also be extended. The current system
already provides a serial REPL, IP-based access with a web interface,
mounted LittleFS transfer, boot logging, and boot-time execution through
\texttt{main.py}. Future work could unify these into a more formal
device-management workflow with remote script deployment, structured
runtime tracing, automated collection of telemetry and images, artifact
export for both obstacle and vision messages, and reproducible
experiment packaging. Such tooling would be especially useful if the
platform is to be used by multiple researchers or integrated into longer
experimental campaigns.

At the system level, the project also opens a path toward broader
hardware generalization. One of the arguments made throughout this work
is that the EdgeTPU is a sufficiently capable low-power inference
primitive to justify close integration with an MCU. A natural next step
is therefore to explore how much of the present design can be ported to
other MCU families, other camera configurations, or future open-source
boards with similar accelerator-centric architectures. If successful,
this would strengthen the broader claim that lightweight
accelerator-plus-MCU designs are a viable class of embedded autonomy
platform rather than a single-board exception.

Taken together, these directions suggest that the present work is less
an endpoint than a foundation. The current implementation establishes
the runtime, sensing, memory, control, and experimentation abstractions
needed for a compact embedded autonomy stack. Future work can now build
on this base to improve viewpoint continuity, flight-platform symmetry,
regulation-aware perception, deployment tooling, embedded learning
workflows, and cross-platform generalization.
